\documentclass[10pt,a4paper]{amsart}
\usepackage[margin=19mm]{geometry}
\usepackage{amsmath,amssymb,mathtools}
\usepackage[scr=rsfs,cal=euler]{mathalfa}
\usepackage{xfrac}
\usepackage[unicode,hidelinks,bookmarks=false]{hyperref}
\newcommand{\ie}{i.e.\ }
\newcommand{\cf}{cf.\ }
\newcommand{\resp}{respectively\ }
\newcommand{\Subsection}{\subsection}
\providecommand{\nobreakdash}{\nobreak\hbox{-}}
\setlength{\emergencystretch}{2em}
\multlinegap0pt
\usepackage{xcolor}
\usepackage{amssymb}
\usepackage{tcolorbox}
\usepackage{xfrac}
\usepackage[normalem]{ulem}
\usepackage{tikz}
\newtheorem{theorem}{Theorem}
\newcommand{\ZZ}{{\mathbb Z}}
\title{Spectra for finite unions of line segments}
\author{Mihail N. Kolountzakis, Ruxi Shi, Sha Wu}
\date{Source: arXiv:2512.19872v1 (CC BY 4.0)}
\begin{document}
\maketitle

\noindent\emph{Source and licence.} Spectra for finite unions of line segments. Version arXiv:2512.19872v1 (CC BY 4.0), distributed by arXiv under the Creative Commons Attribution 4.0 International licence, http://creativecommons.org/licenses/by/4.0/, as recorded in the arXiv licence metadata for this exact version and shown on the abstract page https://arxiv.org/abs/2512.19872v1. Full licence notice: \texttt{licenses/arxiv-2512.19872-CC-BY-4.0.txt}.\par\smallskip

\noindent\textit{Editorial scope.} Counted unit: the article's theorem twoline1-4 (source0.tex lines 609--613). The article's complete original proof of it is delivered with it (source0.tex lines 614--690, one \texttt{proof} environment of 77 lines). No mathematics is written, repaired or re-derived: the statement and the proof are the article's own lines, in the article's own notation, and no retained line is substituted or re-flowed.  No further source-local context is needed: every cross-reference of every retained line resolves inside the two delivered spans.  Nothing was omitted from any retained span, so the package carries no omission record.  No retained line contains a citation, so the wrapper carries no bibliography and no bibliography entry.  No retained line contains a figure environment, an image-inclusion command or drawing code, so no image is delivered and the package declares no asset dependency.  Editorial formatting only, in addition: an AMS \texttt{amsart} preamble that restates the article's own declarations and macros used by the retained lines, namely the environments \texttt{theorem} on separate counters, together with the standard packages those lines need.  The retained environments print in this document as Theorem 1 in order of appearance, whereas the article prints the same items with the numbering of its own full document.  The complete native source of the article as distributed in the arXiv e-print for this exact version is bundled unchanged as \texttt{sources/191516/source0.tex}.
\begin{theorem}\label{twoline1-4}
 Let $T\geq1$ and $m$ be a positive integer. If   $\Lambda =\{\alpha_1, \alpha_2, \cdots ,\alpha_{2m}\}+m\mathbb{Z}$
with   $0=\alpha_1< \alpha_2< \cdots \alpha_{2m}<m $ and  $$
\sum_{\lambda  \in \Lambda }\left|\hat{ \mathbf{1}}_{\left[0, T \right]}\right|^2\left(s-\lambda\right) = 2T $$ for   $s \in \mathbb{R}$, then $ \Lambda=\{0, \alpha\}+\mathbb{Z}$  for some $0<\alpha<1$. Moreover,  we have  $ \Lambda=   \frac{\mathbb{Z}}{2}$ when $T>1$. 
\end{theorem}

\begin{proof}
 Applying the Fourier transform to both sides of the equation 
 $$\sum_{\lambda  \in \Lambda }\left|\hat{ \mathbf{1}}_{\left[0, T \right]}\right|^2\left(s-\lambda\right) = 2T$$
 and the distributional Poisson summation formula, we get
$$
( { \mathbf{1}}_{\left[0, T\right]}* \widetilde{ \mathbf{1}}_{\left[0, T\right]})(x) \cdot \sum_{j=1} e^{2 \pi i \alpha_j x}\cdot \frac{1}{m} \sum_{k \in \mathbb{Z}} \delta_{\frac{k}{m}}=2 T \delta_0.
$$
Since  $\hat{ \mathbf{1}}_{\left[0, T\right]}* \tilde{ \mathbf{1}}_{\left[0, T\right]} > 0$ in $(-T, T)$ and $T\geq 1$, we conclude that
$$
\sum_{j=1}^{2m} e^{2 \pi i \alpha_j\frac{l}{m}}=0\ \ \text{ for any }\ \ l\in\{ \pm 1,  \pm2, \cdots,  \pm  (m-1)\}.
$$  
Let $u_j=e^{2 \pi i \frac{\alpha_j}{m}}$ and  
 denote    by  $p_k\left(u_1, \ldots, u_{2m}\right)$   the $k$-th power sum
$$
p_k\left(u_1, u_2, \ldots, u_{2m}\right)=\sum_{i=1}^{2m} u_i^k=u_1^k+\cdots+u_{2m}^k 
\ \   \text{ for} \ \  k \geqslant 1.
$$
Then we have 
\begin{equation}\label{eq(0.002)}
\begin{cases}
p_1\left(u_1, u_2, \ldots, u_{2m}\right)=0,\\[6pt]
p_2\left(u_1, u_2, \ldots, u_{2m}\right)=0,\\[6pt]
  \cdots  \\%[6pt]
 p_{m-1}\left(u_1, u_2, \ldots, u_{2m}\right)=0. 
 \end{cases}
\end{equation}
The polynomial with roots $u_j$  may be expanded as
$$
P(x):=\prod_{j=1}^{2m}\left(x-u_j\right)=\sum_{j=0}^{2m}(-1)^j e_j x^{2m-j},
$$
where (these are the so-called Newton identities)
\begin{equation*}
\begin{cases} 
e_0 =1, \\[6pt]
e_1 =p_1, \\[6pt]
e_2 =\frac{1}{2}\left(e_1 p_1-p_2\right),\\[6pt]
e_3 =\frac{1}{3}\left(e_2 p_1-e_1 p_2+p_3\right),\\%[6pt]
 \vdots
\end{cases}
\end{equation*}
Here $e_j$ is the $j$-th elementary symmetric function of the numbers $u_1, \ldots, u_{2m}$.

By \eqref{eq(0.002)}, we have $e_1 =e_2 = \cdots =e_{m-1}=0$. This means 
\begin{align*}
P(x) &= x^{2m}+(-1)^{m}e_{m} x^{m}+(-1)^{m+1}e_{m+1} x^{m-1}\\
  &\hspace{3em} +(-1)^{m+2}e_{m+2} x^{m-2}+\cdots+(-1)^{2m}e_{2m}\\
 &= x^{2m}+0\cdot x^{2m-1}+0\cdot x^{2m-2}+\cdots +0\cdot x^{m+1} + (-1)^{m}e_{m}  x^{m} \\
 &\hspace{3em}+(-1)^{m+1}e_{m+1} x^{m-1}+(-1)^{m+2}e_{m+2} x^{m-2}+\cdots+e_{2m}.
\end{align*}
Moreover, $P\left(\frac{1}{x}\right) =\prod_{i=1}^{2m}\left(\frac{1}{x}-u_i\right)=\sum_{j=0}^{2m}(-1)^j e_j x^{j-2m}$ and 
\begin{align*} 
x^{2m}P\left(\frac{1}{x}\right)& = 1+(-1)^{m}e_{m} x^{m} +(-1)^{m+1}e_{m+1} x^{m+1}\\
 &\hspace{3em}+(-1)^{m+2}e_{m+2} x^{m+2}+\cdots+(-1)^{2m}e_{2m}x^{2m}
\\&= e_{2m}x^{2m} - e_{2m-1} x^{2m-2}+\cdots
%+ (-1)^{m+2}e_{m+2} x^{m+2}
+ (-1)^{m+1}e_{m+1}x^{m+1}  \\
&\hspace{3em} + (-1)^{m }e_{m } x^{m}  +0\cdot x^{m-1}+\cdots+ 0\cdot x  +1.
\end{align*}
Note that   since the roots of $P(x)$ are $u_1, u_2,\cdots u_{2m}$ and  the roots of $x^{2m}P(\frac{1}{x})$  are $\overline{u_1}, \overline{u_2},\cdots \overline{u_{2m}}$, it follows that  $$P(x)=Cx^{2m}\overline{P}(\frac{1}{x})$$ for $C = e_{2m}$ and  $u_1, u_2,\cdots u_{2m}$  are  the roots of $$x^{2m}+(-1)^{m}e_m x^m+C=0.$$
Since  $u_1=1$ it follows that the numbers $u_1, \ldots, u_{2m}$ are precisely the numbers  $$\{e^{2\pi i  \frac{j}{m}}: j\in\{0,1, \cdots, m-1\}\}\cup\{e^{2\pi i  \frac{\alpha+j}{m}}: j\in\{0,1, \cdots, m-1\}  \}$$ for some $\alpha\in(0, 1)$. It follows that  the numbers $a_1, \ldots, a_{2m}$ are  exactly the numbers
$$0, \alpha, 1, \alpha+1, \ldots, m-1, \alpha+m-1$$ and therefore 
$ \Lambda=\{0, \alpha\}+\mathbb{Z}$.

Moreover, if $T>1$, then 
$$
\sum_{j=1}^{2m} e^{2 \pi i \alpha_j\frac{l}{m}}=0\ \ \text{ for any }\ \ l\in\{ \pm 1,  \pm2, \cdots,  \pm  m\}.
$$
So
\ifdefined\SMART
$$p_1\left(u_1, u_2, \ldots, u_{2m}\right)=\cdots=p_m\left(u_1, u_2, \ldots, u_{2m}\right)=0$$
\else
$$p_1\left(u_1, u_2, \ldots, u_{2m}\right)=p_2\left(u_1, u_2, \ldots, u_{2m}\right)=\cdots=p_m\left(u_1, u_2, \ldots, u_{2m}\right)=0$$
\fi
and $e_1=e_2=\cdots=e_m=0$. From the proof above, we can obtain $ u_1, u_2,\cdots u_{2m}$  are  the roots of $x^{2m}+C=0$. Since  $u_1=1$, this means that $C=-1$ and the numbers $u_1, \ldots, u_{2m}$ are precisely the numbers $$ \{e^{2\pi i\frac{j}{2m}}: j=0, 1, \ldots, 2m-1\}.$$ It follows that  the numbers $a_1, \ldots, a_{2m}$ are the numbers
$0, \frac12, 1, \ldots, m-\frac12$ and therefore $\Lambda = \frac12\ZZ$.

\end{proof}

\end{document}
